\chapter{2006年福建卷（理）}
\section{单选题}
\begin{problem}
设 \(a\)，\(b\)，\(c\)，\(d\in\R\)，则复数 \((a + b\i)(c + d\i)\) 为实数的充要条件是（\quad）

\choices
    {\(ad - bc = 0\)}
    {\(ac - bd = 0\)}
    {\(ac + bd = 0\)}
    {\(ad + bc = 0\)}
\end{problem}

\begin{answer}
D
\end{answer}

\begin{solution}
由复数乘法得 \((a+b\i)(c+d\i)=(ac-bd)+(ad+bc)\i\). 因为 \(a\)，\(b\)，\(c\)，\(d\in\R\)，复数为实数当且仅当虚部系数为零，即 \(ad+bc=0\)，所以选 D.
\end{solution}

\begin{problem}
在等差数列 \(\{a_{n}\}\) 中，已知 \(a_1 = 2\)，\(a_2 + a_3 = 13\)，则 \(a_4 + a_5 + a_6\) 等于（\quad）

\choices
    {40}
    {42}
    {43}
    {45}
\end{problem}

\begin{answer}
B
\end{answer}

\begin{solution}
设等差数列 \(\{a_n\}\) 的公差为 \(d\)，则 \(a_2=2+d\)，\(a_3=2+2d\). 由题意，\(a_2+a_3=4+3d=13\)，解得 \(d=3\). 因而 \(a_4=11\)，\(a_5=14\)，\(a_6=17\)，所以 \(a_4+a_5+a_6=11+14+17=42\)，选 B.
\end{solution}

\begin{problem}
已知 \(\alpha \in \left(\frac{\pi}{2}, \pi\right)\)，\(\sin \alpha = \frac{3}{5}\)，则 \(\tan \left(\alpha + \frac{\pi}{4}\right)\) 等于（\quad）

\choices
    {\( \frac{1}{7} \)}
    {7}
    {\(-\frac{1}{7}\)}
    {\(-7\)}
\end{problem}

\begin{answer}
A
\end{answer}

\begin{solution}
因为 \(\alpha\in\left(\frac{\pi}{2},\pi\right)\)，所以 \(\cos\alpha<0\). 由 \(\sin\alpha=\frac{3}{5}\) 得 \(\cos\alpha=-\frac{4}{5}\)，从而 \(\tan\alpha=-\frac{3}{4}\). 因此 \(\tan\left(\alpha+\frac{\pi}{4}\right)=\frac{\tan\alpha+1}{1-\tan\alpha}=\frac{-\frac{3}{4}+1}{1+\frac{3}{4}}=\frac{1}{7}\)，所以选 A.
\end{solution}

\begin{problem}
已知全集 \(U = \R\)，且 \(A = \{x \mid |x - 1| > 2\}\)，\(B = \{x \mid x^2 - 6x + 8 < 0\}\)，则 \((\complement_U A) \cap B\) 等于（\quad）

\choices
    {\([-1, 4)\)}
    {\((2,3)\)}
    {\((2,3]\)}
    {\((-1,4)\)}
\end{problem}

\begin{answer}
C
\end{answer}

\begin{solution}
由 \(|x-1|>2\) 得 \(x<-1\) 或 \(x>3\)，所以 \(A=(-\infty,-1)\cup(3,+\infty)\)，从而 \(\complement_UA=[-1,3]\). 又 \(x^2-6x+8=(x-2)(x-4)\)，故 \(B=(2,4)\). 因此 \((\complement_UA)\cap B=[-1,3]\cap(2,4)=(2,3]\)，所以选 C.
\end{solution}

\begin{problem}
已知正方体外接球的体积是 \(\frac{32}{3}\pi\)，那么正方体的棱长等于（\quad）

\choices
    {\(2\sqrt{2}\)}
    {\( \frac{2\sqrt{3}}{3} \)}
    {\( \frac{4\sqrt{2}}{3} \)}
    {\( \frac{4\sqrt{3}}{3} \)}
\end{problem}

\begin{answer}
D
\end{answer}

\begin{solution}
设正方体棱长为 \(s\)，外接球半径为 \(R\). 由球的体积，\(\frac{4}{3}\pi R^3=\frac{32}{3}\pi\)，得 \(R=2\). 正方体的体对角线为 \(s\sqrt{3}\)，且它等于外接球直径 \(2R=4\)，所以 \(s\sqrt{3}=4\)，即 \(s=\frac{4\sqrt{3}}{3}\)，选 D.
\end{solution}

\begin{problem}
在一个口袋中装有 5 个白球和 3 个黑球，这些球除颜色外完全相同，从中摸出 3 个球，至少摸到 2 个黑球的概率等于（\quad）

\choices
    {\(\frac{2}{7}\)}
    {\(\frac{3}{8}\)}
    {\(\frac{3}{4}\)}
    {\(\frac{9}{28}\)}
\end{problem}

\begin{answer}
A
\end{answer}

\begin{solution}
从 \(8\) 个球中任取 \(3\) 个，基本事件总数为 \(\mathrm C_8^3\). 至少摸到 \(2\) 个黑球包括摸到 \(2\) 个黑球和 \(3\) 个黑球两种情况，有利事件数为 \(\mathrm C_3^2\mathrm C_5^1+\mathrm C_3^3=15+1=16\). 所以所求概率为 \(\frac{16}{\mathrm C_8^3}=\frac{16}{56}=\frac{2}{7}\)，选 A.
\end{solution}

\begin{problem}
对于平面 \(\alpha\) 和共面的直线 \(m\)，\(n\)，下列命题中真命题是（\quad）

\choices
    {若 \(m \perp \alpha\)，\(m \perp n\)，则 \(n \parallel \alpha\)}
    {若 \(m \parallel \alpha\)，\(n \parallel \alpha\)，则 \(m \parallel n\)}
    {若 \(m \subset \alpha\)，\(n \parallel \alpha\)，则 \(m \parallel n\)}
    {若 \(m\)，\(n\) 与 \(\alpha\) 所成的角相等，则 \(m \parallel n\)}
\end{problem}

\begin{answer}
C
\end{answer}

\begin{solution}
选项 A 不正确：取 \(\alpha\) 为 \(z=0\) 平面，\(m\) 为 \(z\) 轴，\(n\) 为 \(x\) 轴，则 \(m\perp\alpha\)，\(m\perp n\)，但 \(n\subset\alpha\)，所以 \(n\) 不平行于 \(\alpha\). 选项 B 不正确：在同一个与 \(\alpha\) 平行的平面内取两条相交的直线 \(m\)，\(n\)，则它们都平行于 \(\alpha\)，但彼此不平行. 选项 D 不正确：若 \(m\)，\(n\) 都在 \(\alpha\) 内且不平行，则它们与 \(\alpha\) 所成的角都为 \(0^\circ\)，但 \(m\) 与 \(n\) 不平行. 对于选项 C，\(m\subset\alpha\)，而 \(n\parallel\alpha\)，所以 \(n\) 不与 \(\alpha\) 相交，从而 \(m\) 与 \(n\) 不相交；又因 \(m\)，\(n\) 共面，故 \(m\parallel n\). 所以选 C.
\end{solution}

\begin{problem}
函数 \( y = \log_{2} \frac{x}{x - 1} \)（\(x > 1\)）的反函数是（\quad）

\choices
    {\(y = \frac{2^x}{2^x - 1}\)（\(x > 0\)）}
    {\(y = \frac{2^x}{2^x - 1}\)（\(x < 0\)）}
    {\(y = \frac{2^x - 1}{2^x}\)（\(x > 0\)）}
    {\(y = \frac{2^x - 1}{2^x}\)（\(x < 0\)）}
\end{problem}

\begin{answer}
A
\end{answer}

\begin{solution}
设原函数中的自变量为 \(x\)，函数值为 \(y\). 由 \(y=\log_{2}\frac{x}{x-1}\) 得 \(2^y=\frac{x}{x-1}\)，所以 \(x(2^y-1)=2^y\)，即 \(x=\frac{2^y}{2^y-1}\). 当 \(x>1\) 时，\(\frac{x}{x-1}>1\)，故 \(y>0\). 交换 \(x\)，\(y\)，反函数为 \(y=\frac{2^x}{2^x-1}\)（\(x>0\)），所以选 A.
\end{solution}

\begin{problem}
已知函数 \( f(x) = 2\sin \omega x \)（\(\omega > 0\)）在区间 \(\left[-\frac{\pi}{3}, \frac{\pi}{4}\right]\) 上的最小值是 \(-2\)，则 \(\omega\) 的最小值等于（\quad）

\choices
    {\( \frac{2}{3} \)}
    {\( \frac{3}{2} \)}
    {2}
    {3}
\end{problem}

\begin{answer}
B
\end{answer}

\begin{solution}
函数 \(f(x)\) 取得最小值 \(-2\) 当且仅当 \(\sin(\omega x)=-1\)，即存在 \(k\in\Z\) 使 \(\omega x=-\frac{\pi}{2}+2k\pi\). 当 \(x\in\left[-\frac{\pi}{3},\frac{\pi}{4}\right]\) 时，\(\omega x\in\left[-\frac{\omega\pi}{3},\frac{\omega\pi}{4}\right]\). 要使这个区间包含最接近原点的取值 \(-\frac{\pi}{2}\)，必须有 \(\frac{\omega\pi}{3}\geq\frac{\pi}{2}\)，即 \(\omega\geq\frac{3}{2}\). 当 \(\omega=\frac{3}{2}\) 时，取 \(x=-\frac{\pi}{3}\)，便有 \(\omega x=-\frac{\pi}{2}\)，函数值为 \(-2\). 其他取值 \(-\frac{\pi}{2}+2k\pi\) 的绝对值不小于 \(\frac{3\pi}{2}\)，所需的 \(\omega\) 更大，因此 \(\omega\) 的最小值为 \(\frac{3}{2}\)，选 B.
\end{solution}

\begin{problem}
已知双曲线 \(\frac{x^2}{a^2} - \frac{y^2}{b^2} = 1\)（\(a > 0\)，\(b > 0\)）的右焦点为 \(F\)，若过点 \(F\) 且倾斜角为 \(60^\circ\) 的直线与双曲线的右支有且只有一个交点，则此双曲线离心率的取值范围是（\quad）

\choices
    {\((1,2]\)}
    {\((1,2)\)}
    {\([2, +\infty)\)}
    {\((2, +\infty)\)}
\end{problem}

\begin{answer}
C
\end{answer}

\begin{solution}
设双曲线的焦半距为 \(c\)，离心率为 \(e=\frac{c}{a}>1\)，则 \(c=ae\)，\(b^2=c^2-a^2=a^2(e^2-1)\). 过右焦点 \(F=(c,0)\) 且倾斜角为 \(60^\circ\) 的直线为 \(y=\sqrt{3}(x-c)\). 令 \(X=\frac{x}{a}\)，代入双曲线方程，得到
\[
(e^2-1)X^2-3(X-e)^2=e^2-1
\]
即 \((e^2-4)X^2+6eX+1-4e^2=0\). 当 \(e\ne2\) 时，该方程的判别式为 \(\Delta=16(e^2-1)^2>0\)，两个根化为 \(X_1=\frac{2e+1}{e+2}\)，\(X_2=\frac{2e-1}{2-e}\). 双曲线右支上的点满足 \(X\geq1\). 对 \(1<e<2\)，有 \(X_1-1=\frac{e-1}{e+2}>0\) 且 \(X_2-1=\frac{3(e-1)}{2-e}>0\)，所以有两个右支交点；对 \(e>2\)，有 \(X_1>1\) 且 \(X_2<0\)，所以只有一个右支交点. 当 \(e=2\) 时，代入方程得 \(12X-15=0\)，只有一个根 \(X=\frac{5}{4}>1\)，也只有一个右支交点. 因此离心率的取值范围为 \([2,+\infty)\)，选 C.
\end{solution}

\begin{problem}
已知 \(\left|\overrightarrow{OA}\right|=1\)，\(\left|\overrightarrow{OB}\right|=\sqrt{3}\)，\(\overrightarrow{OA}\cdot\overrightarrow{OB}=0\)，点 \(C\) 在 \( \angle AOB \) 内，且 \( \angle AOC=30^{\circ}\). 设 \( \overrightarrow{OC}=m\overrightarrow{OA}+n\overrightarrow{OB} \)（\(m\)，\(n\in\R\)），则 \( \frac{m}{n} \) 等于（\quad）

\choices
    {\( \frac{1}{3} \)}
    {3}
    {\( \frac{\sqrt{3}}{3} \)}
    {\(\sqrt{3}\)}
\end{problem}

\begin{answer}
B
\end{answer}

\begin{solution}
令 \(\bs e_1=\overrightarrow{OA}\)，\(\bs e_2=\frac{\overrightarrow{OB}}{\sqrt{3}}\)，则 \(\bs e_1\)，\(\bs e_2\) 是互相垂直的单位向量. 由 \(\overrightarrow{OC}=m\overrightarrow{OA}+n\overrightarrow{OB}\)，有 \(\overrightarrow{OC}=m\bs e_1+n\sqrt{3}\bs e_2\). 点 \(C\) 在 \(\angle AOB\) 内，所以 \(m>0\)，\(n>0\)，且相对于 \(\overrightarrow{OA}\) 的正切值为 \(\frac{n\sqrt{3}}{m}\). 因为 \(\angle AOC=30^\circ\)，故 \(\frac{n\sqrt{3}}{m}=\tan30^\circ=\frac{1}{\sqrt{3}}\)，从而 \(m=3n\)，即 \(\frac{m}{n}=3\)，选 B.
\end{solution}

\begin{problem}
对于直角坐标平面内的任意两点 \(A(x_{1}, y_{1})\)，\(B(x_{2}, y_{2})\)，定义它们之间的一种“距离”：\(\|AB\| = |x_{2} - x_{1}| + |y_{2} - y_{1}|\). 给出下列三个命题：

\begin{circlelist}
\item 若点 \(C\) 在线段 \(AB\) 上，则 \( \|AC\| + \|CB\| = \|AB\| \)；
\item 在 \(\triangle ABC\) 中，若 \(\angle C = 90^{\circ}\)，则 \(\|AC\|^{2} + \|CB\|^{2} = \|AB\|^{2}\)；
\item 在 \(\triangle ABC\) 中，\(\|AC\| + \|CB\| > \|AB\|\).
\end{circlelist}

其中真命题的个数为（\quad）

\choices
    {1}
    {2}
    {3}
    {4}
\end{problem}

\begin{answer}
A
\end{answer}

\begin{solution}
对第一个命题，若 \(C=(1-t)A+tB\)，其中 \(0\leq t\leq1\)，则 \(C-A=t(B-A)\)，\(B-C=(1-t)(B-A)\). 因而 \(\|AC\|=t\|AB\|\)，\(\|CB\|=(1-t)\|AB\|\)，所以 \(\|AC\|+\|CB\|=\|AB\|\)，第一个命题为真. 对第二个命题，取 \(A=(0,0)\)，\(C=(1,0)\)，\(B=(1,1)\)，则 \(\angle C=90^\circ\)，但 \(\|AC\|=\|CB\|=1\)，\(\|AB\|=2\)，所以 \(\|AC\|^2+\|CB\|^2=2\ne4=\|AB\|^2\)，第二个命题为假. 对第三个命题，取 \(A=(0,0)\)，\(B=(2,2)\)，\(C=(1,0)\)，三点不共线，但 \(\|AC\|=1\)，\(\|CB\|=3\)，\(\|AB\|=4\)，等号成立而不是严格大于，第三个命题为假. 因此真命题有 \(1\) 个，选 A.
\end{solution}

\section{填空题}
\begin{problem}
\(\left(x^{2} - \frac{1}{x}\right)^{5}\) 展开式中 \(x^{4}\) 的系数是\fillinblank{}（用数字作答）.
\end{problem}

\begin{answer}
\(10\)
\end{answer}

\begin{solution}
在展开式中，若选取 \(r\) 个因子 \(-\frac{1}{x}\)，其余 \(5-r\) 个因子选取 \(x^2\)，所得项的系数为 \(\mathrm C_5^r(-1)^r\)，次数为 \(2(5-r)-r=10-3r\). 要得到 \(x^4\)，需满足 \(10-3r=4\)，解得 \(r=2\). 因而 \(x^4\) 的系数为 \(\mathrm C_5^2(-1)^2=10\).
\end{solution}

\begin{problem}
已知直线 \(x - y - 1 = 0\) 与抛物线 \(y = ax^2\) 相切，则 \(a = \)\fillinblank{}.
\end{problem}

\begin{answer}
\(\frac{1}{4}\)
\end{answer}

\begin{solution}
直线 \(x-y-1=0\) 可写为 \(y=x-1\). 联立直线与抛物线，得 \(ax^2=x-1\)，即 \(ax^2-x+1=0\). 相切时该关于 \(x\) 的二次方程有且只有一个实根，因此判别式满足 \(1-4a=0\)，解得 \(a=\frac{1}{4}\).
\end{solution}

\begin{problem}
一个均匀小正方体的六个面中，三个面上标以数 \(0\)，两个面上标以数 \(1\)，一个面上标以数 \(2\). 将这个小正方体抛掷 \(2\text{ 次}\)，则向上的数之积的数学期望是\fillinblank{}.
\end{problem}

\begin{answer}
\(\frac{4}{9}\)
\end{answer}

\begin{solution}
设一次抛掷向上的数为随机变量 \(X\). 由六个面等可能向上，\(P(X=0)=\frac{3}{6}\)，\(P(X=1)=\frac{2}{6}\)，\(P(X=2)=\frac{1}{6}\)，所以 \(E(X)=0\cdot\frac{3}{6}+1\cdot\frac{2}{6}+2\cdot\frac{1}{6}=\frac{2}{3}\). 两次抛掷相互独立，设两次向上的数为 \(X_1\)，\(X_2\)，则 \(E(X_1X_2)=E(X_1)E(X_2)=\left(\frac{2}{3}\right)^2=\frac{4}{9}\).
\end{solution}

\begin{problem}
如图，连结 \(\triangle ABC\) 的各边中点得到一个新的 \(\triangle A_{1}B_{1}C_{1}\)，又连结 \(\triangle A_{1}B_{1}C_{1}\) 的各边中点得到一个新的 \(\triangle A_{2}B_{2}C_{2}\)，如此无限继续下去，得到一系列三角形：\(\triangle ABC\)，\(\triangle A_{1}B_{1}C_{1}\)，\(\triangle A_{2}B_{2}C_{2}\)，\(\cdots\)，这一系列三角形趋向于一个点 \(M\). 已知 \(A(0,0)\)，\(B(3,0)\)，\(C(2,2)\)，则点 \(M\) 的坐标是\fillinblank{}.
\bitmapfigure[max width=0.55\linewidth]{img/2006/fujian_science/q16_fig1.png}
\end{problem}

\begin{answer}
\(\left(\frac{5}{3},\frac{2}{3}\right)\)
\end{answer}

\begin{solution}
记初始三角形为 \(A_0B_0C_0=ABC\). 每次取边中点时，
\(\overrightarrow{OA_{k+1}}=\frac{1}{2}\left(\overrightarrow{OB_k}+\overrightarrow{OC_k}\right)\)，\(\overrightarrow{OB_{k+1}}=\frac{1}{2}\left(\overrightarrow{OC_k}+\overrightarrow{OA_k}\right)\)，\(\overrightarrow{OC_{k+1}}=\frac{1}{2}\left(\overrightarrow{OA_k}+\overrightarrow{OB_k}\right)\).
将这三个等式相加，得到
\[
\overrightarrow{OA_{k+1}}+\overrightarrow{OB_{k+1}}+\overrightarrow{OC_{k+1}}=\overrightarrow{OA_k}+\overrightarrow{OB_k}+\overrightarrow{OC_k}
\]
所以每个三角形的重心都相同. 另一方面，每次取边中点后，三角形的边长变为原来的一半，故其直径趋于 \(0\)，三个顶点都趋于同一点 \(M\). 因此 \(M\) 就是初始三角形的重心，故
\[
M=\left(\frac{0+3+2}{3},\frac{0+0+2}{3}\right)=\left(\frac{5}{3},\frac{2}{3}\right)
\]
\end{solution}

\section{解答题}
\begin{problem}
已知函数 \(f(x) = \sin^2 x + \sqrt{3}\sin x\cos x + 2\cos^2 x\)，\(x \in \R\).

\begin{enumerate}
\item 求函数 \( f(x) \) 的最小正周期和单调增区间；
\item 函数 \( f(x) \) 的图象可以由函数 \( y = \sin 2x \)（\(x \in \R\)）的图象经过怎样的变换得到？
\end{enumerate}
\end{problem}

\begin{solution}
\begin{enumerate}
\item 利用倍角公式，得
\[
f(x)=\frac{1-\cos 2x}{2}+\frac{\sqrt{3}}{2}\sin 2x+1+\cos 2x
=\frac{3}{2}+\frac{\sqrt{3}}{2}\sin 2x+\frac{1}{2}\cos 2x
=\frac{3}{2}+\sin\left(2x+\frac{\pi}{6}\right)
\]
因为正弦函数的最小正周期为 \(2\pi\)，所以 \(f(x)\) 的最小正周期为 \(\pi\). 又 \(f'(x)=2\cos\left(2x+\frac{\pi}{6}\right)\)，故当
\[
-\frac{\pi}{2}+2k\pi<2x+\frac{\pi}{6}<\frac{\pi}{2}+2k\pi
\]
时，\(f'(x)>0\)，即 \(k\pi-\frac{\pi}{3}<x<k\pi+\frac{\pi}{6}\)，其中 \(k\in\Z\). 因此，\(f(x)\) 的单调增区间为 \(\left[k\pi-\frac{\pi}{3},k\pi+\frac{\pi}{6}\right]\)，其中 \(k\in\Z\).
\item 由上式可得 \(f(x)=\sin\left(2\left(x+\frac{\pi}{12}\right)\right)+\frac{3}{2}\). 因此，将 \(y=\sin 2x\) 的图象向左平移 \(\frac{\pi}{12}\) 个单位长度，再向上平移 \(\frac{3}{2}\) 个单位长度，即可得到函数 \(f(x)\) 的图象.
\end{enumerate}
\end{solution}

\begin{problem}
如图，四面体 \(ABCD\) 中，\(O\)，\(E\) 分别是 \(BD\)，\(BC\) 的中点，\(CA = CB = CD = BD = 2\)，\(AB = AD = \sqrt{2}\).
\begin{enumerate}
\item 求证：\(AO \perp\) 平面 \(BCD\)；
\item 求异面直线 \(AB\) 与 \(CD\) 所成角的大小；
\item 求点 \(E\) 到平面 \(ACD\) 的距离.
\end{enumerate}

\bitmapfigure[max width=0.45\linewidth]{img/2006/fujian_science/q18_fig1.png}
\end{problem}

\begin{solution}
\begin{enumerate}
\item 在三角形 \(ABD\) 中，\(AB=AD=\sqrt{2}\)，\(BD=2\)，且 \(AB^{2}+AD^{2}=BD^{2}\)，所以 \(\angle BAD=90^{\circ}\). 由于 \(O\) 是直角三角形 \(ABD\) 的斜边 \(BD\) 的中点，所以 \(OA=OB=OD=1\). 又 \(AB=AD\)，点 \(A\) 在 \(BD\) 的垂直平分面上，故 \(AO\perp BD\).

在三角形 \(BCD\) 中，\(BC=CD=BD=2\)，所以三角形 \(BCD\) 是等边三角形，\(CO\perp BD\)，且 \(CO=\sqrt{3}\). 在三角形 \(AOC\) 中，\(OA^{2}+OC^{2}=1+3=4=AC^{2}\)，所以 \(AO\perp CO\). 直线 \(BD\) 与 \(CO\) 是平面 \(BCD\) 内相交的两条直线，故 \(AO\perp\) 平面 \(BCD\).
\item 取 \(O\) 为原点，令平面 \(BCD\) 为 \(xOy\) 平面，建立直角坐标系，使 \(B=(1,0,0)\)，\(D=(-1,0,0)\)，\(C=(0,\sqrt{3},0)\)，\(A=(0,0,1)\). 于是 \(\overrightarrow{AB}=(1,0,-1)\)，\(\overrightarrow{CD}=(-1,-\sqrt{3},0)\)，从而 \(\left|\overrightarrow{AB}\cdot\overrightarrow{CD}\right|=1\)，\(\left|\overrightarrow{AB}\right|=\sqrt{2}\)，\(\left|\overrightarrow{CD}\right|=2\).
设异面直线 \(AB\) 与 \(CD\) 所成的角为 \(\theta\)，则
\[
\cos\theta=\frac{\left|\overrightarrow{AB}\cdot\overrightarrow{CD}\right|}{\left|\overrightarrow{AB}\right|\left|\overrightarrow{CD}\right|}=\frac{1}{2\sqrt{2}}=\frac{\sqrt{2}}{4}
\]
所以 \(\theta=\arccos\frac{\sqrt{2}}{4}\).
\item 由 \(A=(0,0,1)\)，\(C=(0,\sqrt{3},0)\)，\(D=(-1,0,0)\)，可取平面 \(ACD\) 的一个法向量为
\[
\overrightarrow{AD}\times\overrightarrow{AC}=(\sqrt{3},-1,-\sqrt{3})
\]
故平面 \(ACD\) 的方程为 \(\sqrt{3}x-y-\sqrt{3}z+\sqrt{3}=0\). 又 \(E\) 是 \(BC\) 的中点，所以 \(E=\left(\frac{1}{2},\frac{\sqrt{3}}{2},0\right)\). 因此点 \(E\) 到平面 \(ACD\) 的距离为
\[
\frac{\left|\sqrt{3}\cdot\frac{1}{2}-\frac{\sqrt{3}}{2}-\sqrt{3}\cdot0+\sqrt{3}\right|}{\sqrt{3+1+3}}=\frac{\sqrt{3}}{\sqrt{7}}=\frac{\sqrt{21}}{7}
\]
\end{enumerate}
\end{solution}

\begin{problem}
统计表明，某种型号的汽车在匀速行驶中每小时的耗油量 \(y\)（升）关于行驶速度 \(x\)（千米/小时）的函数解析式可以表示为：\(y=\frac{1}{128000}x^3-\frac{3}{80}x+8\)（\(0<x\leqslant120\)）. 已知甲、乙两地相距 \(100\text{ 千米}\).

\begin{enumerate}
\item 当汽车以 \(40\text{ 千米/小时}\) 的速度匀速行驶时，从甲地到乙地要耗油多少升？
\item 当汽车以多大的速度匀速行驶时，从甲地到乙地耗油最少？ 最少为多少升？
\end{enumerate}
\end{problem}

\begin{solution}
\begin{enumerate}
\item 当汽车以 \(40\text{ 千米/小时}\) 的速度行驶时，每小时耗油量为
\[
y=\frac{40^{3}}{128000}-\frac{3}{80}\times40+8=7
\]
行驶 \(100\text{ 千米}\) 所需时间为 \(\frac{100}{40}=\frac{5}{2}\) 小时，所以耗油量为 \(7\times\frac{5}{2}=\frac{35}{2}\text{ 升}\).
\item 当速度为 \(x\text{ 千米/小时}\) 时，行驶 \(100\text{ 千米}\) 所需时间为 \(\frac{100}{x}\) 小时，耗油量为
\[
H(x)=\frac{100}{x}\left(\frac{x^{3}}{128000}-\frac{3}{80}x+8\right)=\frac{x^{2}}{1280}-\frac{15}{4}+\frac{800}{x}
\]
其中 \(0<x\leqslant120\). 求导得
\[
H'(x)=\frac{x}{640}-\frac{800}{x^{2}}=\frac{x^{3}-512000}{640x^{2}}=\frac{x^{3}-80^{3}}{640x^{2}}
\]
因为 \(x>0\)，所以当 \(0<x<80\) 时 \(H'(x)<0\)，当 \(80<x\leqslant120\) 时 \(H'(x)>0\). 因此 \(H(x)\) 在 \(x=80\) 时取得最小值，且
\[
H(80)=\frac{80^{2}}{1280}-\frac{15}{4}+\frac{800}{80}=\frac{45}{4}\text{ 升}
\]
所以汽车以 \(80\text{ 千米/小时}\) 的速度匀速行驶时耗油最少，最少为 \(\frac{45}{4}\text{ 升}\).
\end{enumerate}
\end{solution}

\begin{problem}
已知椭圆 \(\frac{x^2}{2} + y^2 = 1\) 的左焦点为 \(F\)，\(O\) 为坐标原点.
\begin{enumerate}
\item 求过点 \(O\)，\(F\)，并且与椭圆的左准线 \(l\) 相切的圆的方程；
\item 设过点 \(F\) 且不与坐标轴垂直的直线交椭圆于 \(A\)，\(B\) 两点，线段 \(AB\) 的垂直平分线与 \(x\) 轴交于点 \(G\)，求点 \(G\) 横坐标的取值范围.
\end{enumerate}

\bitmapfigure[max width=0.42\linewidth]{img/2006/fujian_science/q20_fig1.png}
\end{problem}

\begin{solution}
\begin{enumerate}
\item 由椭圆方程可知 \(a^{2}=2\)，\(b^{2}=1\)，所以 \(c=\sqrt{a^{2}-b^{2}}=1\). 因此左焦点 \(F=(-1,0)\)，左准线 \(l\) 的方程为 \(x=-\frac{a^{2}}{c}=-2\).

设所求圆的圆心为 \(M\). 因为圆经过 \(O\) 和 \(F\)，所以 \(M\) 在 \(OF\) 的垂直平分线上，可设 \(M=\left(-\frac{1}{2},t\right)\). 圆的半径满足
\[
r=|MO|=\sqrt{\frac{1}{4}+t^{2}}
\]
圆与直线 \(l\) 相切，所以半径等于圆心到 \(l\) 的距离，即 \(r=\left|-\frac{1}{2}-(-2)\right|=\frac{3}{2}\). 于是
\[
\frac{1}{4}+t^{2}=\frac{9}{4}
\]
所以 \(t=\pm\sqrt{2}\).
所以所求圆的方程为
\[
\left(x+\frac{1}{2}\right)^{2}+\left(y-\sqrt{2}\right)^{2}=\frac{9}{4}
\]
或
\[
\left(x+\frac{1}{2}\right)^{2}+\left(y+\sqrt{2}\right)^{2}=\frac{9}{4}
\]
\item 设直线 \(AB\) 的方程为 \(y=k(x+1)\). 由于该直线不与坐标轴垂直，所以 \(k\neq0\)，且它不是竖直直线. 代入椭圆方程，得
\[
\frac{x^{2}}{2}+k^{2}(x+1)^{2}=1
\Longleftrightarrow
(1+2k^{2})x^{2}+4k^{2}x+2k^{2}-2=0
\]
设交点 \(A\)，\(B\) 的横坐标分别为 \(x_{1}\)，\(x_{2}\)，则由韦达定理
\[
x_{1}+x_{2}=-\frac{4k^{2}}{1+2k^{2}}
\]
设线段 \(AB\) 的中点为 \(N(x_{0},y_{0})\)，于是 \(x_{0}=\frac{x_{1}+x_{2}}{2}=-\frac{2k^{2}}{1+2k^{2}}\)，\(y_{0}=k(x_{0}+1)=\frac{k}{1+2k^{2}}\).
因为 \(AB\) 的斜率为 \(k\)，所以其垂直平分线的斜率为 \(-\frac{1}{k}\)，方程为 \(y-y_{0}=-\frac{1}{k}(x-x_{0})\). 令 \(y=0\)，得到点 \(G\) 的横坐标
\[
x_G=x_{0}+ky_{0}=-\frac{2k^{2}}{1+2k^{2}}+\frac{k^{2}}{1+2k^{2}}=-\frac{k^{2}}{1+2k^{2}}
\]
令 \(u=k^{2}>0\)，则
\[
x_G=-\frac{u}{1+2u}=-\frac{1}{2}+\frac{1}{2(1+2u)}
\]
由于 \(u>0\)，故 \(-\frac{1}{2}<x_G<0\). 所以点 \(G\) 横坐标的取值范围为 \(\left(-\frac{1}{2},0\right)\).
\end{enumerate}
\end{solution}

\begin{problem}
已知函数 \( f(x) = -x^2 + 8x \)，\( g(x) = 6\ln x + m \).

\begin{enumerate}
\item 求 \( f(x) \) 在区间 \( \left[t, t + 1\right] \) 上的最大值 \( h(t) \)；
\item 是否存在实数 \(m\)，使得 \(y = f(x)\) 的图象与 \(y = g(x)\) 的图象有且只有三个不同的交点？若存在，求出 \(m\) 的取值范围；若不存在，说明理由.
\end{enumerate}
\end{problem}

\begin{solution}
\begin{enumerate}
\item 因为 \(f(x)=-(x-4)^{2}+16\)，所以 \(f(x)\) 在 \((-\infty,4]\) 上单调递增，在 \([4,+\infty)\) 上单调递减. 于是当 \(t<3\) 时，区间 \([t,t+1]\) 在 \(4\) 的左侧，最大值为 \(f(t+1)\)；当 \(3\leqslant t\leqslant4\) 时，区间包含 \(4\)，最大值为 \(f(4)\)；当 \(t>4\) 时，区间在 \(4\) 的右侧，最大值为 \(f(t)\). 因此
\[
h(t)=
\begin{cases}
-t^{2}+6t+7, & t<3 \\
16, & 3\leqslant t\leqslant4 \\
-t^{2}+8t, & t>4
\end{cases}
\]
\item 交点的横坐标必须满足 \(x>0\) 且
\[
-x^{2}+8x=6\ln x+m
\]
令 \(\varphi(x)=-x^{2}+8x-6\ln x\)，则交点个数等于方程 \(\varphi(x)=m\) 在 \(x>0\) 上的根的个数. 求导得
\[
\varphi'(x)=-2x+8-\frac{6}{x}=-\frac{2(x-1)(x-3)}{x}
\]
所以 \(\varphi(x)\) 在 \((0,1)\) 上单调递减，在 \((1,3)\) 上单调递增，在 \((3,+\infty)\) 上单调递减. 又 \(\lim_{x\to0^{+}}\varphi(x)=+\infty\)，\(\varphi(1)=7\)，\(\varphi(3)=15-6\ln3\)，\(\lim_{x\to+\infty}\varphi(x)=-\infty\).
令 \(q(x)=\frac{1}{2}\left(x-\frac{1}{x}\right)-\ln x\)，则 \(q(1)=0\)，且 \(q'(x)=\frac{(x-1)^2}{2x^2}>0\)（\(x>1\)），所以 \(\ln3<\frac{1}{2}\left(3-\frac{1}{3}\right)=\frac{4}{3}\)，从而 \(15-6\ln3>7\).
因此水平直线 \(y=m\) 与三段图象各有一个交点当且仅当
\[
7<m<15-6\ln3
\]
因此存在这样的实数 \(m\)，且其取值范围为 \(\left(7,15-6\ln3\right)\).
\end{enumerate}
\end{solution}

\begin{problem}
已知数列 \(\{a_{n}\}\) 满足 \(a_1 = 1\)，\(a_{n+1} = 2a_n + 1\)（\(n \in \N^*\)）.

\begin{enumerate}
\item 求数列 \(\{a_{n}\}\) 的通项公式；
\item 若数列 \(\{b_{n}\}\) 满足 \(4^{b_1 - 1}4^{b_2 - 1}4^{b_3 - 1}\cdots 4^{b_n - 1} = (a_n + 1)^{b_n}\)（\(n \in \N^*\)），证明 \(\{b_n\}\) 是等差数列；
\item 证明：\(\frac{n}{2} - \frac{1}{3} < \frac{a_1}{a_2} + \frac{a_2}{a_3} + \ldots + \frac{a_n}{a_{n+1}} < \frac{n}{2}\)（\(n \in \N^*\)）.
\end{enumerate}
\end{problem}

\begin{solution}
\begin{enumerate}
\item 由递推式得 \(a_{n+1}+1=2(a_n+1)\)，而 \(a_1+1=2\)，所以 \(\{a_n+1\}\) 是首项为 \(2\)、公比为 \(2\) 的等比数列，从而 \(a_n+1=2^n\)，即
\[
a_n=2^n-1
\]
\item 由第（1）问，原等式可写成
\[
4^{b_1-1}4^{b_2-1}\cdots4^{b_n-1}=(2^n)^{b_n}
\]
令 \(S_n=b_1+b_2+\cdots+b_n\)，则上式两边都化为以 \(2\) 为底的幂，故
\[
2S_n-2n=nb_n
\]
分别取 \(n\) 和 \(n-1\)，相减得
\[
2b_n-2=nb_n-(n-1)b_{n-1}
\Longleftrightarrow
(n-2)b_n-(n-1)b_{n-1}+2=0
\]
令 \(n=2\)，得 \(b_1=2\). 再令 \(d=b_2-2\). 当 \(n=3\) 时，上式给出 \(b_3=2b_2-2=2+2d\). 若对某个 \(k\geqslant3\) 有 \(b_k=2+(k-1)d\)，则在上式中取 \(n=k+1\)，得
\[
b_{k+1}=\frac{kb_k-2}{k-1}=\frac{k\left[2+(k-1)d\right]-2}{k-1}=2+kd
\]
因此由数学归纳法，\(b_n=2+(n-1)d\) 对一切 \(n\in\N^*\) 成立，故 \(\{b_n\}\) 是等差数列.
\item 由 \(a_n=2^n-1\)，有
\[
\frac{a_k}{a_{k+1}}=\frac{2^k-1}{2^{k+1}-1}=\frac{1}{2}-\frac{1}{2(2^{k+1}-1)}<\frac{1}{2}
\]
对 \(k\geqslant1\)，因为
\[
2(2^{k+1}-1)-3\cdot2^k=2^k-2\geqslant0
\]
所以
\[
0<\frac{1}{2(2^{k+1}-1)}\leqslant\frac{1}{3\cdot2^k}
\]
于是
\[
0<\sum_{k=1}^{n}\frac{1}{2(2^{k+1}-1)}\leqslant\sum_{k=1}^{n}\frac{1}{3\cdot2^k}<\frac{1}{3}
\]
将上式代入
\[
\sum_{k=1}^{n}\frac{a_k}{a_{k+1}}=\frac{n}{2}-\sum_{k=1}^{n}\frac{1}{2(2^{k+1}-1)}
\]
可得
\[
\frac{n}{2}-\frac{1}{3}<\frac{a_1}{a_2}+\frac{a_2}{a_3}+\cdots+\frac{a_n}{a_{n+1}}<\frac{n}{2}
\]
\end{enumerate}
\end{solution}
